% Source: https://artofproblemsolving.com/community/c3h1877656_2018_national_sprint_30_wsolution_sketch
% Stable identity: aops-topic-1877656-q01
Both triangles are $3$-$4$-$5$ right triangles, with right angles at $C$ and $E$. Put $A=(0,0)$, $D=(5,0)$, $C=(4,0)$, $B=(4,3)$. A reflection followed by a translation reverses orientation, so $E$ lies on the other side of $AD$ from $B$: from $AE=3$, $DE=4$ we get $E=\left(\dfrac{9}{5},-\dfrac{12}{5}\right)$.

\textbf{The axis passes through midpoints.} Let $g$ be the reflection in $m$ followed by a translation along $m$. For any point $P$, the reflected point $P_1$ is the mirror image of $P$, so the midpoint of $PP_1$ lies on $m$; translating $P_1$ along $m$ moves that midpoint along $m$ (by half the translation). Hence the midpoint of $P$ and $g(P)$ always lies on $m$.

Applying this to $A\mapsto D$, $B\mapsto A$, $C\mapsto E$, line $m$ passes through
\[
K=\text{midpoint of }AD=\left(\dfrac{5}{2},0\right),\quad M=\text{midpoint of }AB=\left(2,\dfrac{3}{2}\right),\quad L=\text{midpoint of }CE=\left(\dfrac{29}{10},-\dfrac{6}{5}\right).
\]
(These are collinear: both $MK$ and $KL$ have slope $-3$, so $m$ is the line $y=-3x+\dfrac{15}{2}$. One checks that reflecting in this line and then translating by $\left(\dfrac{1}{2},-\dfrac{3}{2}\right)$, a vector parallel to $m$, sends $A,B,C$ to $D,A,E$.)

\begin{asy}
usepackage("amsmath");
size(6cm);
defaultpen(fontsize(10pt));
pair A=(0,0),D=(5,0),C=(4,0),B=(4,3),E2=(1.8,-2.4);
pair M=(A+B)/2,K=(A+D)/2,L=(C+E2)/2;
fill(A--M--K--cycle,paleblue);
draw(A--B--C--cycle,linewidth(1));draw(D--A--E2--cycle,linewidth(1));
pair u=unit(K-M);
draw((M-1.3*u)--(L+1.1*u),red+dashed+linewidth(0.9));
draw(C--E2,gray+dotted+linewidth(0.8));
dot(A);dot(B);dot(C);dot(D);dot(E2);dot(M,red);dot(K,red);dot(L,red);
label("$A$",A,W);label("$B$",B,N);label("$C$",C,NE);label("$D$",D,E);label("$E$",E2,S);
label("$M$",M,NW);label("$K$",K+(-0.33,-0.33));label("$L$",L+(-0.36,-0.06));
label("$m$",L+1.1*u,E,red);
label("\mbox{}",(min(currentpicture,user=true)+max(currentpicture,user=true))/2,black); // keeps the text after the figure black when LaTeX typesets the labels inline
\end{asy}

So $m$ meets side $AB$ at its midpoint $M$ and side $AC$ at $K$ with $AK=\dfrac{5}{2}$, cutting off triangle $AMK$ (shaded). Since triangles $AMK$ and $ABC$ share the angle at $A$,
\[
[AMK]=\dfrac{AM}{AB}\cdot\dfrac{AK}{AC}\cdot[ABC]=\dfrac{1}{2}\cdot\dfrac{5}{8}\cdot6=\dfrac{15}{8}.
\]
The other region has area $6-\dfrac{15}{8}=\dfrac{33}{8}$, so the ratio is
\[
\dfrac{\dfrac{15}{8}}{\dfrac{33}{8}}=\boxed{\dfrac{5}{11}}.
\]
